% Unregistered research auxiliary; not a new main-gate conclusion.
\paragraph{Fixed-parameter negative-$K$ obstruction (proved auxiliary statement).}
Let $k$ be an algebraically closed field of characteristic $p>0$, let
$(R,\mathfrak m)$ be a noetherian local $k$-algebra with residue field $k$,
and let $g:Z\to\operatorname{Spec}R$ be a proper morphism from a regular
noetherian scheme which is an isomorphism away from the closed point.
Put $E=(Z\times_R k)_{\mathrm{red}}$.  Let $P$ be a fixed smooth proper
$k$-scheme; in the intended application $P$ is a connected curve.  All
$K$-groups below are Bass-completed nonconnective groups of perfect
complexes.  We have a natural isomorphism
\begin{equation}
 \operatorname{coker}\!\left(
 K_0(Z\times_k P)[1/p]\longrightarrow K_0(E\times_k P)[1/p]
 \right)
 \ \simeq\ K_{-1}(\operatorname{Spec}R\times_k P)[1/p].
 \tag{FF.1}
\end{equation}
Thus an actual class on the fixed product $E\times_k P$ lifts to an actual
class on $Z\times_k P$, after inverting $p$, if and only if its image
under the connecting map in (FF.1) is zero.  This is a statement about
$K_0$ classes, not about a compatible system of chosen vector bundles.

\emph{Proof.}
We write $B=\operatorname{Spec}R\times_k P$ and $Z_P=Z\times_k P$.
Because $P$ is of finite presentation over $k$, the scheme $B$ is
noetherian.  The morphism $Z_P\to B$ is proper and is an isomorphism
off the closed subscheme $P=\operatorname{Spec}k\times_k P$.
It therefore gives an abstract blow-up square.  Its infinitesimal
centres and exceptional schemes are
\[
 B_n=\operatorname{Spec}(R/\mathfrak m^n)\times_k P,
 \qquad
 F_n=Z\times_R\operatorname{Spec}(R/\mathfrak m^n)\times_k P.
\]
We use \citet[Theorem A and (1.6), author pp.~2--3]
{rootKerzStrunkTamme2018} on this square.  In particular we keep the term
$K_0(B)$; the product $B$ is not being treated as a local scheme.

The immersions $P\hookrightarrow B_n$ and $E\times_kP\hookrightarrow F_n$
are nilpotent.  For the second assertion, the schemes have the same
underlying closed set, and their ideals are nilpotent on the noetherian
scheme $F_n$; products with the smooth $k$-scheme $P$ preserve the
reduced scheme $E$ because $k$ is perfect.  On every affine open,
\citet[Theorem A, p.~29]{rootGeisserHesselholt2011} gives a
$p$-primary relative group in every integer degree.  Applying
\citet[Theorem 8.1, p.~367]{D_ThomasonTrobaugh1990} to a finite affine
cover globalizes the equivalence after inverting $p$.  We obtain,
compatibly with all transition maps,
\begin{equation}
 K_i(B_n)[1/p]\simeq K_i(P)[1/p],\qquad
 K_i(F_n)[1/p]\simeq K_i(E\times_kP)[1/p]
 \quad(i\in\mathbf Z).
 \tag{FF.2}
\end{equation}
We do not claim a bound for the $p$-exponents uniform in $n$ or $i$.

Tensoring the exact pro-group sequence of KST levelwise with
$\mathbf Z[1/p]$ is exact.  One way to see the only categorical point
needed here is to represent a morphism of pro-abelian groups on a common
indexing category: kernels and cokernels are then represented levelwise,
and the localization functor is exact on abelian groups.  Using (FF.2),
all terms in the relevant segment become constant pro-groups.  The
constant embedding of abelian groups into pro-abelian groups is exact
and reflects exactness.  Hence we have an ordinary exact sequence
\begin{align*}
 K_0(B)[1/p]&\longrightarrow
 K_0(Z_P)[1/p]\oplus K_0(P)[1/p]
 \longrightarrow K_0(E\times_kP)[1/p]\\
 &\xrightarrow{\ \partial_P\ }K_{-1}(B)[1/p]
 \longrightarrow K_{-1}(Z_P)[1/p]\oplus K_{-1}(P)[1/p].
\end{align*}
Both $Z_P$ and $P$ are regular: the first is smooth over the regular
scheme $Z$.  Their negative groups vanish by
\citet[Proposition 6.8(b), pp.~362--363]{D_ThomasonTrobaugh1990}.
Finally, the image of $K_0(P)$ in $K_0(E\times_kP)$ already lies in the
image of $K_0(Z_P)$, via the actual projection $Z_P\to P$.  Exactness
therefore gives (FF.1).  No inverse limit has been formed in this proof.
\qed

\paragraph{Products and the proposed tetrahedral coefficient.}
The same argument gives the natural exact sequence in every integer
degree.  In particular, for $P=\operatorname{Spec}k$ it gives
\[
 K_{-2}(E)[1/p]\xrightarrow[\sim]{\ \partial\ }
 K_{-3}(R)[1/p].
\]
The sequences are compatible with external products by classes on the
fixed parameter scheme.  Indeed, each relative $K$-theory fibre is a
module over the base $K$-theory spectrum; the connecting map of a fibre
sequence of modules is linear, with the usual graded convention.
The nonconnective module structure is described in
\citet[Section 6.5, pp.~360--361]{D_ThomasonTrobaugh1990}.
Thus, if $\kappa\in K_{-2}(E)[1/p]$ and $\xi\in K_2(P)[1/p]$, then
\begin{equation}
 \partial_P(\kappa\boxtimes\xi)
   = (\partial\kappa)\boxtimes\xi
   \quad\hbox{in }K_{-1}(\operatorname{Spec}R\times_kP)[1/p].
 \tag{FF.3}
\end{equation}
The degree of $\xi$ is even, so exchanging the displayed factors does
not introduce a sign.

\paragraph{Henselian point parameters and the remaining family obstruction.}
If $R$ is henselian, then
\begin{equation}
 K_0(Z)[1/p]\xrightarrow{\ \sim\ }K_0(E)[1/p].
 \tag{FF.4}
\end{equation}
Indeed, \citet[Theorem 3.7 and Remark (ii), author p.~8]
{frDrinfeld2004Tate} proves $K_{-1}(R)=0$ for every commutative henselian
local ring.  This gives surjectivity by (FF.1).  For injectivity, we
use the preceding exact sequence with $P=\operatorname{Spec}k$.
The homomorphism $K_0(R)\to K_0(k)$ is the rank isomorphism
$\mathbf Z\to\mathbf Z$, so the image of $K_0(R)[1/p]$ in
$K_0(Z)[1/p]\oplus\mathbf Z[1/p]$ is the diagonal graph of the unit
class.  Its intersection with $K_0(Z)[1/p]\oplus0$ is zero.
This proves injectivity directly from exactness.  In particular every
individual proposed coefficient class $\kappa\cdot u$ in $K_0(E)[1/p]$
lifts.  The scheme $\operatorname{Spec}R\times_kP$ for a positive-dimensional
proper parameter $P$ is not henselian local; Drinfeld's theorem does
not dispose of its group in (FF.1).

For the tetrahedral regular modification constructed separately in
(TG.1)--(TG.5), the whole-diagram calculation (TS.1)--(TS.4) identifies
the dual-complex coefficient in the indicated ordinary $K$-group.
Formula (FF.3) gives the exact obstruction for lifting its entire
proper-curve family.  Neither the nonvanishing of $\partial\kappa$ nor the existence
of a proper-curve $K_2$ relation proves that its product in (FF.3) is
nonzero or zero.  We assert neither outcome.  In particular (FF.1)
does not imply family surjectivity merely from nilinvariance on every
finite thickening.  The pointwise isomorphism (FF.4) does not settle a
rational ordinary-Chow gate.  A separate endpoint argument may produce
an actual family with the same two endpoint restrictions without
lifting the originally prescribed class over the whole parameter;
that conclusion would not assert the vanishing of (FF.3).


\paragraph{Dated scope update after C.3bu.}
The preceding obstruction formulas remain valid in their full
$\mathbf Z[1/p]$ scope.  In the stronger geometric hypotheses of
C.3bn, with $d\leq2$ and $P$ a smooth proper curve, the later
C.3bu theorem now proves rational surjectivity on the entire
prescribed $K_0$ family and hence
$K_{-1}(\operatorname{Spec}R\times_kP)_{\mathbf Q}=0$.
Thus the earlier retained whole-family caution is superseded in
exactly that rational, fixed-local-base, proper-curve scope.
The historical attempted coefficient detector is preserved above.
No general $\mathbf Z[1/p]$ vanishing, higher-dimensional parameter
statement, affine-family compactification, or transcendental
nonlocal-base assertion follows.  The intrinsic component
Albanese-kernel-family problem of C.3bt remains unchanged.
